%% kapa: exam=1 round=34 year=2023 period=1 form=A subject=economics subject_ko=경제학원론 seq=32 number=72
%% source: Q-Net 감정평가사 기출문제 (cst00309) artlSeq=5211789 file=raw/1008/5211789/2023년도 제34회 감정평가사 1차시험 1교시.pdf
%% conversion: kordoc 4.24.51 -> tools/gen1.py (auto); review=proofread-sonnet; flags=none
\begin{question}{72}
폐쇄경제의 IS-LM 모형에서 화폐시장 균형조건이 $\frac{M}{P} = L(r, Y - T)$일 때, 조세삭감이 미치는 효과로 옳은 것을 모두 고른 것은? (단, 초기는 균형상태, IS곡선은 우하향, LM곡선은 우상향하며, $M$은 통화량, $P$는 물가, $r$은 이자율, $Y$는 국민소득, $T$는 조세이다.)
\begin{framedtext}
\begin{tabular}{@{}p{0.45\linewidth}p{0.45\linewidth}@{}}
ㄱ. IS곡선 우측이동 & ㄴ. LM곡선 우측이동 \\
ㄷ. 통화수요 감소 & ㄹ. 이자율 상승 \\
\end{tabular}
\end{framedtext}
\begin{choices}
  \item ㄱ, ㄴ
  \item ㄱ, ㄷ
  \item ㄱ, ㄹ
  \item ㄴ, ㄷ
  \item ㄴ, ㄹ
\end{choices}
\end{question}
\answer{3}
